\documentclass[12pt,a4paper]{article}
\usepackage[utf8]{inputenc}
\usepackage{amsmath}
\usepackage{graphicx}
\usepackage{float}
\usepackage{listings}
\usepackage{xcolor}
\usepackage{geometry}
\usepackage{tabularx}
\usepackage{hyperref}
\geometry{a4paper, margin=1in}

\definecolor{codegreen}{rgb}{0,0.6,0}
\definecolor{codegray}{rgb}{0.5,0.5,0.5}
\definecolor{codepurple}{rgb}{0.58,0,0.82}
\definecolor{backcolour}{rgb}{0.95,0.95,0.92}

\lstdefinestyle{mystyle}{
    backgroundcolor=\color{backcolour},   
    commentstyle=\color{codegreen},
    keywordstyle=\color{magenta},
    numberstyle=\tiny\color{codegray},
    stringstyle=\color{codepurple},
    basicstyle=\ttfamily\footnotesize,
    breakatwhitespace=false,         
    breaklines=true,                 
    captionpos=b,                    
    keepspaces=true,                 
    numbers=left,                    
    numbersep=5pt,                  
    showspaces=false,                
    showstringspaces=false,
    showtabs=false,                  
    tabsize=2,
    columns=fullflexible,
    postbreak=\mbox{\textcolor{red}{$\hookrightarrow$}\space}
}
\lstset{style=mystyle}

\title{\textbf{LAB 1 - DELIVERABLES} \\ Sequential Multiplier Using Repeated Addition}
\author{\textbf{Name:} R SARANG \\ \textbf{Roll Number:} EC23I2015 \\ \textbf{Course:} VLSI Testing Practice (EC5034)}
\date{30th August 2026}

\begin{document}

\maketitle
\tableofcontents
\newpage

\section*{Aim}
To design and verify a 4-bit unsigned sequential multiplier using repeated addition, consisting of one adder and an FSM-based control path.

\vspace{0.5em}
\noindent \textbf{GitHub Repository:} \url{https://github.com/rsarang14/VLSI-Testing-Practice}

\section{Algorithm}
The sequential multiplier computes the product of two 4-bit unsigned numbers (multiplicand $A$ and multiplier $B$) using repeated addition. 
The algorithm initializes an 8-bit accumulator to zero. In each computation step, it adds the multiplicand to the accumulator and decrements the multiplier by 1. This process is repeated until the multiplier reaches zero. 

\begin{figure}[H]
    \centering
    \includegraphics[width=0.6\textwidth]{handwritten-1.png}
    \caption{Algorithm Pseudocode}
\end{figure}

\section{Datapath Elements}

\begin{figure}[H]
    \centering
    \includegraphics[width=0.7\textwidth]{handwritten-2.png}
    \caption{Datapath Element Usage}
\end{figure}
\begin{itemize}
    \item \textbf{Accumulated Product Register:} 8-bit register to hold the running sum.
    \item \textbf{Working Multiplicand Register:} 4-bit register to store operand A.
    \item \textbf{Loop Counter:} 4-bit down-counter to store operand B and decrement each cycle.
    \item \textbf{Adder:} 8-bit adder to add the working multiplicand to the accumulated product.
    \item \textbf{Zero Detector:} Combinational logic to check if the loop counter is zero.
\end{itemize}

\section{Datapath Block Diagram}
\begin{figure}[H]
    \centering
    \includegraphics[width=0.7\textwidth]{handwritten-3.png}
    \caption{Complete Datapath Block Diagram}
\end{figure}

\section{Control Signals}
Control signals are generated by the FSM (Controller) and used to orchestrate the operations in the Datapath.
\begin{table}[H]
\centering
\begin{tabularx}{\textwidth}{|l|l|X|}
\hline
\textbf{Signal} & \textbf{Direction} & \textbf{Function} \\ \hline
\texttt{load\_inputs} & Controller $\rightarrow$ Datapath & Loads the input operands ($A$ into multiplicand reg, $B$ into counter). \\ \hline
\texttt{clear\_prod} & Controller $\rightarrow$ Datapath & Clears the 8-bit product register (accumulator) to 0. \\ \hline
\texttt{enable\_prod} & Controller $\rightarrow$ Datapath & Enables product update (addition of multiplicand to accumulator). \\ \hline
\texttt{enable\_cnt\_dec} & Controller $\rightarrow$ Datapath & Enables loop counter operation (decrements the counter by 1). \\ \hline
\texttt{mult\_complete} & Controller $\rightarrow$ Output & Indicates completion of the multiplication. \\ \hline
\end{tabularx}
\end{table}

\section{Status Signals}
Status signals are generated by the Datapath and provided to the FSM (Controller) to make transition decisions.
\begin{table}[H]
\centering
\begin{tabularx}{\textwidth}{|l|l|X|}
\hline
\textbf{Signal} & \textbf{Direction} & \textbf{Function} \\ \hline
\texttt{loops\_zero\_flag} & Datapath $\rightarrow$ Controller & Indicates that the counter has reached zero. \\ \hline
\end{tabularx}
\end{table}

\section{Timing Diagram}
\begin{figure}[H]
    \centering
    \includegraphics[width=0.85\textwidth]{handwritten-4.png}
    \caption{Timing Diagram for Control and Data Signals}
\end{figure}

\section{FSM Design}
\begin{figure}[H]
    \centering
    \includegraphics[width=0.75\textwidth]{handwritten-5.png}
    \caption{FSM State Table and State Diagram}
\end{figure}

\begin{table}[H]
\centering
\begin{tabular}{|c|c|c|c|l|}
\hline
\textbf{Current State} & \textbf{start} & \textbf{loops\_zero\_flag} & \textbf{Next State} & \textbf{Outputs Asserted} \\ \hline
IDLE (00) & 0 & X & IDLE (00) & None \\ \hline
IDLE (00) & 1 & X & LOAD (01) & None \\ \hline
LOAD (01) & X & X & PROG (10) & \texttt{load\_inputs}=1, \texttt{clear\_prod}=1 \\ \hline
PROG (10) & X & 0 & PROG (10) & \texttt{enable\_prod}=1, \texttt{enable\_cnt\_dec}=1 \\ \hline
PROG (10) & X & 1 & DONE (11) & None \\ \hline
DONE (11) & 1 & X & DONE (11) & \texttt{mult\_complete}=1 \\ \hline
DONE (11) & 0 & X & IDLE (00) & \texttt{mult\_complete}=1 \\ \hline
\end{tabular}
\end{table}

\section{Integrate \& Verify}
\subsection{Verilog RTL}
The Verilog RTL design consists of the datapath, controller, and top-level modules.

\subsubsection{Datapath (\texttt{datapath.v})}
\lstinputlisting[language=Verilog]{datapath.v}

\subsubsection{Controller (\texttt{controller.v})}
\lstinputlisting[language=Verilog]{controller.v}

\subsubsection{Top Module (\texttt{top.v})}
\lstinputlisting[language=Verilog]{top.v}

\subsection{Testbench (\texttt{tb\_top.v})}
\lstinputlisting[language=Verilog]{tb\_top.v}

\subsection{Test-Case Results}
\begin{table}[H]
\centering
\begin{tabular}{|c|c|c|c|c|l|}
\hline
\textbf{Test Case} & \textbf{A} & \textbf{B} & \textbf{Expected Product} & \textbf{Actual Product} & \textbf{Result} \\ \hline
Corner (0x0) & 0 & 0 & 0 & 0 & PASS \\ \hline
Corner (0xN) & 0 & 5 & 0 & 0 & PASS \\ \hline
Corner (Nx0) & 7 & 0 & 0 & 0 & PASS \\ \hline
Corner (1x1) & 1 & 1 & 1 & 1 & PASS \\ \hline
Corner (1x15) & 1 & 15 & 15 & 15 & PASS \\ \hline
Corner (15x1) & 15 & 1 & 15 & 15 & PASS \\ \hline
Maximum Case & 15 & 15 & 225 & 225 & PASS \\ \hline
Commutativity 1 & 3 & 12 & 36 & 36 & PASS \\ \hline
Commutativity 2 & 12 & 3 & 36 & 36 & PASS \\ \hline
\end{tabular}
\end{table}

\subsection{Cycle-Count Verification}
\textbf{Expected total cycles:} $B + 2$ (Consists of 1 LOAD cycle + B COMPUTE cycles + 1 DONE cycle)

\begin{table}[H]
\centering
\begin{tabular}{|c|c|c|c|c|c|}
\hline
\textbf{A} & \textbf{B} & \textbf{Expected Product} & \textbf{Expected Cycles} & \textbf{Actual Cycles} & \textbf{Result} \\ \hline
0 & 0 & 0 & 2 & 2 & PASS \\ \hline
3 & 4 & 12 & 6 & 6 & PASS \\ \hline
3 & 12 & 36 & 14 & 14 & PASS \\ \hline
12 & 3 & 36 & 5 & 5 & PASS \\ \hline
15 & 15 & 225 & 17 & 17 & PASS \\ \hline
\end{tabular}
\end{table}

\subsection{3 $\times$ 12 vs 12 $\times$ 3 Observation}
Both $3 \times 12$ and $12 \times 3$ result in the same product ($36$) due to the commutative property of multiplication. However, their execution times differ significantly in a sequential multiplier implemented using repeated addition.
\begin{itemize}
    \item For $3 \times 12$ ($A=3, B=12$), the multiplier adds $3$ to the accumulator twelve times. This requires $12$ COMPUTE cycles, resulting in a total of $14$ cycles ($12+2$).
    \item For $12 \times 3$ ($A=12, B=3$), the multiplier adds $12$ to the accumulator three times. This requires only $3$ COMPUTE cycles, resulting in a total of $5$ cycles ($3+2$).
\end{itemize}
\textbf{Conclusion:} To minimize execution time in a repeated-addition sequential multiplier, the smaller operand should always be chosen as the multiplier ($B$).

\section*{Inference}
The sequential multiplier was successfully designed by separating the logic into a datapath and a control path (FSM). The FSM effectively orchestrates the operations in the datapath, correctly managing the state transitions from IDLE to LOAD, to PROG (computation), and finally to DONE. Using a self-checking testbench, the design was validated across a comprehensive set of test cases including corner cases and worst-case scenarios, confirming its correctness and functional completeness.

\section*{Conclusion}
In this lab, a 4-bit unsigned sequential multiplier was designed and verified using Verilog HDL. The separation of datapath and controller simplified the design process and debugging. Simulation results confirm that the FSM properly controls the repeated addition algorithm, producing accurate results for all test cases and corner conditions. The cycle-count analysis successfully demonstrated the performance dependency on the choice of the multiplier operand.

\end{document}
